\documentclass[11pt]{article}
\usepackage[a4paper,margin=27mm]{geometry}
\usepackage{amsmath,amssymb,amsthm,mathtools,xcolor,booktabs,array}
\newcommand{\PROVED}{\textcolor{green!40!black}{\textbf{PROVED}}}
\newcommand{\OPEN}{\textcolor{red!70!black}{\textbf{OPEN}}}
\newcommand{\AUDIT}{\textcolor{orange!70!black}{\textbf{AUDIT}}}
\newtheorem{theorem}{Theorem}
\newtheorem{lemma}{Lemma}
\newtheorem{proposition}{Proposition}
\newtheorem{corollary}{Corollary}
\title{Compactness audit of a hypothetical URGT$_1$ infractor (v126)}
\date{14 August 2026}

\begin{document}
\maketitle

\section{The spectral probability of an infractor}

Use the notation of v124 and put
\begin{equation}
 W_N(\tau):=\sum_{k\ge1}|\Phi_{N,k}(\tau)|^2.              \tag{1}
\end{equation}
Assume that $N_\nu\to\infty$ and that $e_\nu\in E_{N_\nu}$ satisfy
\begin{equation}
 \sum_{k\ge1}\|L_{N_\nu,k}e_\nu\|^2=1,
 \qquad
 \mathcal E_{\Gamma T,N_\nu}(e_\nu)<1.                  \tag{2}
\end{equation}
Define the output spectral probability
\begin{equation}
 d\mu_\nu(\tau)
 :=W_{N_\nu}(\tau)|\widehat e_\nu(\tau)|^2
       \frac{d\tau}{2\pi}.                               \tag{3}
\end{equation}
Plancherel and (2) give $\mu_\nu(\mathbb R)=1$.

\begin{lemma}[An infractor cannot escape to high frequency --- \PROVED]
The family $(\mu_\nu)$ is tight.  More precisely,
\begin{equation}
 \mu_\nu\{|\tau|\ge R\}
 \le\frac1{\inf_{|\tau|\ge R}g_\Gamma(\tau)}             \tag{4}
\end{equation}
for every sufficiently large $R$.
\end{lemma}

\begin{proof}
The Tate term is nonnegative, so (2) implies
\[
 \int_{\mathbb R}g_\Gamma(\tau)\,d\mu_\nu(\tau)<1.
\]
The multiplier $g_\Gamma$ is even, increasing in $|\tau|$, and tends to
$+\infty$.  Markov's inequality gives (4), whose right side tends to zero.
\end{proof}

Thus every infractor has a weakly convergent subsequence.  In particular,
the proposed mean-zero mode cannot evade the argument by moving its output
to frequencies of order $1/\delta_N$: Gamma would make its energy exceed
one.

\section{The all-depth compact arithmetic profile}

For each fixed $k$, the same product/divisibility calculation as v125
gives
\begin{equation}
 \frac{\log N}{N}\Phi_{N,k}(\tau)
 \longrightarrow
 \frac1{(k-1)!}\frac1{\tau^2+\frac14}                     \tag{5}
\end{equation}
locally uniformly in real $\tau$.  Indeed,
\begin{equation}
 \sum_{m\le N}\frac{\Lambda_k(m)}{\sqrt m}\,m^{-i\tau}
 =\frac{N^{1/2-i\tau}(\log N)^{k-1}}
       {(k-1)!(\frac12-i\tau)},(1+o_{k,R}(1)),            \tag{6}
\end{equation}
and the integral-label return is lower order.

The uniform Witt majorant used in v109 is precisely what is needed to
interchange the $N$-limit and the sum over depths.  Written at the
polynomial rather than scalar-norm level, the required statement is
\begin{equation}
 \boxed{
 \left(\frac{\log N}{N}\right)^2W_N(\tau)
 \longrightarrow
 C_{\rm W}\frac1{(\tau^2+\frac14)^2},
 \qquad
 C_{\rm W}:=\sum_{k\ge1}\frac1{((k-1)!)^2},}              \tag{7}
\end{equation}
locally uniformly in $\tau$.

\paragraph{Status of (7) --- \AUDIT.}
Fixed-depth convergence is proved by (5)--(6).  The scalar all-depth
majorant of v109 does not, by itself, state local uniform convergence of
the polynomials.  To use (7) without a gap one must prove the strengthened
tail estimate
\begin{equation}
 \lim_{K\to\infty}\limsup_{N\to\infty}
 \sup_{|\tau|\le R}
 \left(\frac{\log N}{N}\right)^2
 \sum_{k>K}|\Phi_{N,k}(\tau)|^2=0                         \tag{8}
\end{equation}
for every fixed $R$.  Equation (8), and not fixed-depth PNT, is the exact
uniformity gate in the compactness proof.

\section{The only two possible tight limiting modes}

\begin{lemma}[Shrinking-support classification --- \AUDIT]
Assume (7).  Let $h_\nu$ be supported in intervals of lengths tending to
zero, and suppose the probabilities
\[
 \frac{W_{N_\nu}(\tau)|\widehat h_\nu(\tau)|^2\,d\tau}
      {\int W_{N_\nu}|\widehat h_\nu|^2}
\]
are tight.  Every nonzero weak limit has density
\begin{equation}
 d\mu_*(\tau)=
 \frac{|c_0+c_1i\tau|^2}{Z(\tau^2+\frac14)^2}\,d\tau,
 \qquad (c_0,c_1)\ne(0,0).                                \tag{9}
\end{equation}
\end{lemma}

\begin{proof}[Candidate argument and exact missing bound]
Translate the shrinking source intervals to the origin.  After scalar
normalization on one fixed compact frequency interval carrying positive
limiting mass, the Paley--Wiener functions $\widehat h_\nu$ have
exponential types tending to zero.  The local $L^2$ compactness theorem
for Paley--Wiener spaces and a diagonal extraction give a locally uniform
limit $H$, entire of exponential type zero.  Tightness and (7) give
\[
 \int_{\mathbb R}\frac{|H(\tau)|^2}
      {(\tau^2+\frac14)^2}\,d\tau<\infty.
\]
If the normalized sources are bounded in one fixed distributional order,
Paley--Wiener--Schwartz identifies $H$ as the transform of a distribution
supported at the origin, hence as a polynomial.  The displayed weighted
integrability then forces $\deg H\le1$, giving (9).

Tightness of the \emph{output} measures does not yet prove that uniform
distributional-order bound for the normalized \emph{sources}.  Entire
functions of exponential type zero need not be polynomials without such a
growth bound.  Therefore the classification is retained as an audit gate,
not used as an unconditional theorem.
\end{proof}

\begin{proposition}[Gamma excludes both limiting modes --- \PROVED]
For the constant mode $c_1=0$,
\begin{equation}
 \int g_\Gamma\,d\mu_*>1+\frac1{45}.                     \tag{10}
\end{equation}
For the linear/mean-zero mode $c_0=0$, already the first Gamma channel
satisfies
\begin{equation}
 \int g_0\,d\mu_*=3.                                      \tag{11}
\end{equation}
For arbitrary $(c_0,c_1)$ the odd cross term integrates to zero, so the
Gamma expectation is the convex combination of the two diagonal
expectations and is strictly larger than one.
\end{proposition}

\begin{proof}
Equation (10) is Theorem 4 of v125.  Put $b=\frac12$.  For (11),
\begin{align*}
 \int_{\mathbb R}\frac{\tau^2}{(\tau^2+b^2)^2}\,d\tau
 &=\frac\pi{2b},\\
 \int_{\mathbb R}
 \frac{4\tau^4}{(\tau^2+b^2)^3}\,d\tau
 &=\frac{3\pi}{2b}.
\end{align*}
Their ratio is three.  Since $g_\Gamma\ge g_0$, later Gamma channels can
only increase it.  Finally, all weights are even and the mixed term in
$|c_0+c_1i\tau|^2$ is odd, proving the last assertion.
\end{proof}

\begin{theorem}[No asymptotic infractor, conditional compactness theorem
--- \PROVED]
If (8) and the distributional-order bound in Lemma 2 hold, there is no
sequence satisfying (2).
\end{theorem}

\begin{proof}
Lemma 1 gives a weakly convergent subsequence of the probabilities (3).
Equations (7)--(8) and Lemma 2 identify its limit with (9).
The bound in (2), lower semicontinuity, and nonnegativity of the Tate term
give
\[
                         \int g_\Gamma\,d\mu_*\le1.
\]
Proposition 2 gives the strict opposite inequality.  This contradiction
excludes (2).
\end{proof}

\section{Ledger}

The limiting contradiction is complete once its two compactness gates are
available.  There are three logically separate remaining tasks before
URGT$_1$ can be marked proved:
\begin{enumerate}
 \item prove the polynomial all-depth tail estimate (8);
 \item derive the uniform distributional-order bound in Lemma 2 from the
       output-energy hypothesis, rather than assume it;
 \item certify the finitely many $N$ not covered by the asymptotic
       contradiction.  A failure at one fixed $N$ does not automatically
       produce a sequence $N_\nu\to\infty$.
\end{enumerate}

\begin{center}
\begin{tabular}{>{\raggedright\arraybackslash}p{.57\linewidth}
                >{\centering\arraybackslash}p{.14\linewidth}
                >{\raggedright\arraybackslash}p{.19\linewidth}}
\toprule
Statement & Status & Location \\
\midrule
Spectral tightness of every infractor & \PROVED & (3)--(4) \\
Fixed-depth compact leakage profile & \PROVED & (5)--(6) \\
All-depth polynomial tail gate & \OPEN & (8) \\
Classification of tight shrinking-cell limits & \AUDIT & Lemma 2 \\
Constant limit excluded by Gamma margin & \PROVED & (10), v125 \\
Mean-zero linear limit excluded with ratio $3$ & \PROVED & (11) \\
No sequence $N_\nu\to\infty$ of infractors & conditional & Theorem 3 \\
Finite-$N$ operator audit & \OPEN & after asymptotic gate \\
URGT$_1$, $G$, row (d) & \OPEN & compactness gates plus finite audit \\
\bottomrule
\end{tabular}
\end{center}

\end{document}
